\section{Discussão: relações para investigar conjuntamente}

A análise relacional complementa os rankings ao identificar quais origens e produtos finais compõem seus contrastes. A presença como fornecedor responde a demandas padronizadas; os volumes ambientais combinam estrutura produtiva, intensidade emissora e demanda observada. A concentração acrescenta como o volume se reparte, enquanto o backbone destaca relações sob uma regra local. Nenhuma dessas dimensões substitui as demais, e reunir indicadores não constitui confirmação independente de um mesmo resultado.

Os perfis de energia e construção sugerem diferentes escalas de coordenação. Uma origem associada a muitos produtos finais motiva investigar mudanças nos processos e na energia utilizada, com repercussão contábil em diversas cadeias. Um destino que reúne emissões de outras origens motiva examinar conjuntamente materiais, especificações e processos fornecedores. As frentes tecnológicas e de demanda discutidas por \textcite{rissman2020} contextualizam essa leitura. Contudo, a matriz não identifica contratos, compradores imediatos ou capacidade de intervenção dos agentes representados.

O vínculo energia--construção ilustra uma conexão entre essas perspectivas, sem justificar concentrar toda a ação nessa relação. Seus pesos relativos diferem nas duas pontas e as demais contrapartes continuam relevantes em conjunto. Uma política direcionada a uma única atividade poderia deixar de considerar relações que compõem a atribuição de outros produtos. O diagnóstico serve para formular perguntas sobre conjuntos de setores; avaliar instrumentos exige cenários explícitos, custos, possibilidades de substituição e respostas econômicas que este estudo não estima.

A aproximação ambiental condiciona os contrastes quantitativos apresentados. A interpolação supõe evolução linear entre dois anos, sem observação física de 2015. Há ainda limitações de preços e de alocação setorial documentadas na auditoria. Um fator comum positivo reescala volumes, preservando participações, concentração e seleção pelo filtro; alterações diferentes entre setores podem modificar composições e backbone. A coordenada de presença como fornecedor depende apenas da MIP, mas isso não valida o eixo ambiental da comparação.

O filtro também impõe limites próprios à leitura. As caudas expressam desproporção sob seu modelo nulo, não incerteza dos coeficientes ou efeito econômico de uma ligação. A sensibilidade a $\alpha$ informa o comportamento da seleção, mas não resolve a validade ambiental. Por isso, a interpretação combina rede completa, cobertura e contrapartes, evitando tratar ausência no backbone como irrelevância. O aproveitamento em uma aplicação regional exigirá conferir suas fronteiras e dados; os resultados nacionais não se transferem automaticamente.
